% !TeX root = main.tex
{\fontsize{12}{14}\selectfont\tenantselection}
\subsection*{Normal lighting and local controls}
The lighting proposal starts with 92 PIL PLATO 600 MultiPower 3117955 panels on the 600 mm ceiling grid: sixty 22 W, twenty-eight 28 W and four 38 W, totalling 2256 W. Their nominal fluxes are 2843.45, 3555.30 and 4685.55 lm respectively. Separate circuits serve the enclosed rooms and distributed main-office zones, with a maintained mean illuminance target of 500 lx adopted for the design. With this arrangement, the connected lighting density is 7.301 W/m$^2$ using the adopted 309 m$^2$ office area.

Seven passive manual switches S-A--S-G provide local control of their corresponding circuits. These switches must have a minimum rating of 20 AX, an LED switching capacity of at least 600 W and make/break/inrush compatibility with the selected drivers. Since their phases differ, S-A/B/C use separate enclosures. No borrowed neutral is allowed. Each route passes through its switch and includes a 2 m DB rise plus 1.5 m down and 1.5 m back at the switch, giving an assumed vertical allowance of 5 m per lighting radial. Passive controls add no operating power.
\begin{center}\begin{tabular}{lrrrrl}\toprule
Circuit & Lamps & W & Phase & RCBO & Control\\\midrule

LC-A & 18 & 402 & L1 & B10 & S-A \\
LC-B & 17 & 392 & L2 & B10 & S-B \\
LC-C & 18 & 396 & L3 & B10 & S-C \\
LC-D & 11 & 242 & L3 & B10 & S-D \\
LC-E & 14 & 392 & L2 & B10 & S-E \\
LC-F & 10 & 280 & L1 & B10 & S-F \\
LC-G & 4 & 152 & L3 & B10 & S-G \\

\bottomrule\end{tabular}\end{center}
On the plan, the A/B/C numeric labels correspond to the LED IDs and circuits LC-A/B/C. Similarly, D01--D11 correspond to PN01--PN11, E82--E96 excluding E90 to the northeast LEDs, F01--F10 to UE01--UE10, and G01--G04 to SE01--SE04. Power settings follow the common device schedule. Every luminaire and control has one circuit assignment.
% !TeX root = main.tex
\subsection*{Lighting calculation}
To assess lighting performance, the calculation uses manufacturer LDT intensities given in cd/klm, with each profile scaled by its selected flux: 2843.45 lm at 22 W, 3555.30 lm at 28 W and 4685.55 lm at 38 W. Illuminance is then calculated at each workplane grid cell, with spacing at most 0.2 m, as
\[
E=0.80\sum_j\frac{I_j(C,\gamma)\cos\gamma}{d_j^2}.
\]
A luminaire-to-workplane height of 1.9 m is adopted, corresponding to a 2.7 m ceiling and a 0.8 m workplane. Meanwhile, the maintenance factor of 0.80 assumes periodic cleaning and replacement. Room boundaries use the adopted plan envelopes; each room receives light from its own fittings. Only direct light is included in this calculation.
\begin{center}\begin{tabular}{lrrrr}\toprule
Zone & Lamps & Area m$^2$ & Mean lx & Min/mean\\\midrule
Open office & 67 & 231.87 & 592.9 & 0.136\\
Pantry & 11 & 29.54 & 546.3 & 0.390\\
Upper east room & 10 & 33.68 & 588.5 & 0.267\\
South east room & 4 & 13.07 & 662.6 & 0.175\\
\bottomrule\end{tabular}\end{center}
On this basis, all four mean illuminances exceed the adopted 500 lx target. To provide mounting-height tolerance, the four south-east panels use the 38 W setting. At 2.2 m above the workplane, the respective means are 573.0, 505.8, 545.0 and 589.2 lx; the south-east room retains a 17.8\% margin.

When interpreting these results, the minimum/mean ratios include perimeter cells. Overall, the assessment covers mean direct illuminance; a detailed task-area check would also account for furniture, reflections and glare.

\subsection*{Small power and fixed equipment}
Small-power provision includes nine four-socket floor boxes, three twin wall sockets and the fixed connection FS01. Each floor-box/socket assembly is individually 13 A fused. Accordingly, the demand calculation uses a 13 A continuous capacity per assembly; FS01 supplies WH01, a STIEBEL ELTRON SHC10 GB Eltron 235232 undersink water heater. Its 1.4 kW at 220 V gives 6.364 A at unity power factor; the selected protection is B10 with a 10 A double-pole FCU. Meanwhile, building central AC is supplied as central plant and not duplicated in this office DB.

P1--P4 are A1 rings with fused floor-box connections. In comparison, P5/P6/P8 are A2 radials, while P7 is a fixed radial. Both P6 and P8 serve the PANTRY identified on the source plan, with FB09 and SO01 connected to P6. P8 is a separate socket group for convenient isolation. Kitchen sockets are separated from ordinary office circuits under EMSD6E(3)(a).
\begin{center}\begin{tabular}{llllr}\toprule
Circuit & Points & Phase & Topology & Protection\\\midrule

P1 & FB01, FB02 & L1 & A1 ring & B32 \\
P2 & FB03, FB04 & L1 & A1 ring & B32 \\
P3 & FB05, FB06 & L2 & A1 ring & B32 \\
P4 & FB07, FB08 & L2 & A1 ring & B32 \\
P5 & SO02 & L1 & A2 radial & B32 \\
P6 & FB09, SO01 & L3 & A2 radial & B32 \\
P7 & FS01/WH01 & L3 & Fixed radial & B10 \\
P8 & SO03 & L3 & A2 radial & B32 \\

\bottomrule\end{tabular}\end{center}
In terms of coverage, the defined service envelopes satisfy the A1 maximum of 100 m$^2$ and the A2 maximum of 50 m$^2$. Coverage is assessed by service envelope, allowing overlap between circuits. Ring L/N/CPC return to their own circuit terminals. Main 6/10 mm$^2$ cables terminate in accessible junction/FCU assemblies rated for those conductor sizes. At each point, the local connections use 4 mm$^2$ L/N/CPC, with an assumed length of 0.5 m per point.
\clearpage
% !TeX root = main.tex
\subsection*{FCU and containment details}
Every FB01--FB09 assembly shares one 13 A FCU upstream of all four socket ways. Each twin SO01--SO03 likewise shares one 13 A FCU; FS01 has its own 10 A double-pole FCU and accessible earthed socket for the selected WH01 heater; retain its manufacturer-supplied cord and plug. Ring/main L/N/CPC connections must terminate in an accessible junction assembly rated for 6/10 mm$^2$ cables. At each connection, the FCU is positioned at the tap before the 4 mm$^2$ local spur, with its neutral switched together with the line and its PE kept continuous. With 13 A continuous capacity per socket assembly and the selected 6.364 A heater load as the demand basis, any added points require recalculation.

To maintain the adopted grouping factors, the containment is arranged in the following banks: lighting LC-A--G (seven circuits, $C_g=0.54$); A: P1/P2 (four ring legs, 0.65); B: P3/P4 (four ring legs, 0.65); C: P5/P6/P8 (three radials, 0.70); P7 alone (1.00); incoming feeder alone (1.00). Co-located plan paths show routing corridors, not a shared enclosure. Accordingly, separate Method B runs must provide horizontal clear space between groups exceeding twice the largest adjacent overall cable diameter, in accordance with EMSD A5(3), Note 2; where enclosure dimensions govern, the enclosures must be spaced at least as conservatively and the cable separation must be confirmed. Thermal separation is provided by cable clearance, rather than a partition alone. Within each circuit, L/N/CPC stays together. Additionally, contact with thermal insulation must be avoided. If these conditions cannot be achieved, the circuits must be regrouped and resized before installation.

\subsection*{Protection and board arrangement}
As shown in the complete schematic, the arrangement includes C63 four-pole protection at MS-O1, 4-core 35 mm$^2$ plus CPC16, a 100 A four-pole DB isolator, independent L1/L2/L3/N/PE bars and 15 separate 1P+N Type A 30 mA RCBOs. Every load neutral passes through its own RCBO;  none are shared downstream. Meanwhile, PE bypasses all switches. Additionally, the circuits and phases are to be labelled, and the arrangement must allow each final circuit to be isolated without interrupting the others.

Device selection must ensure that the proposed B-curve final protection is compatible with driver inrush and the breaking-capacity/energy limits established in Task 3. Individual RCBOs limit the effect of a residual-current fault to its own circuit; however, full short-circuit selectivity requires manufacturer coordination data.
\subsection*{Emergency lighting assumption}
Landlord emergency lighting and exit signage use independent essential supplies, separated from normal lighting under EMSD 6B(4). These feeds bypass tenant normal-lighting switches and RCBO load sides. Emergency-system loads are therefore outside this office normal-supply DB.
